\begin{examplebox}
	\textbf{\textcolor{CExample}{例 1（基础）}}

	\textbf{\textcolor{CExample}{题目：}}
	已知四边形 $ABCD$ 内接于圆 $O$，$\angle A=80^\circ$，$\angle B=100^\circ$，求 $\angle C$ 和 $\angle D$。

	\textbf{\textcolor{CExample}{解题步骤：}}
	\begin{description}[style=nextline, leftmargin=2.8em, labelsep=0.5em, itemsep=0.45em]
		\item[\textbf{第一步（识别对角）：}] 在圆内接四边形中，$\angle A$ 与 $\angle C$ 为对角，$\angle B$ 与 $\angle D$ 为对角。
			\par\textbf{理由：} 四边形的对角指不相邻的两个内角。
		\item[\textbf{第二步（应用互补）：}] 由对角互补得
			\[
				\begin{aligned}
					\angle C &= 180^\circ - \angle A \\
					&= 180^\circ - 80^\circ \\
					&= 100^\circ .
			\end{aligned}
			\]
			\par\textbf{理由：} $\angle A$ 与 $\angle C$ 互补。
		\item[\textbf{第三步（求另一角）：}] 同理，
			\[
				\begin{aligned}
					\angle D &= 180^\circ - \angle B \\
					&= 180^\circ - 100^\circ \\
					&= 80^\circ .
			\end{aligned}
			\]
			\par\textbf{理由：} $\angle B$ 与 $\angle D$ 互补。
	\end{description}

	\textbf{\textcolor{CExample}{关键结论：}}
	$\angle C=100^\circ$，$\angle D=80^\circ$。

	\medskip
	\textbf{\textcolor{CExample}{例 2（稍变形）}}

	\textbf{\textcolor{CExample}{题目：}}
	圆内接四边形 $ABCD$ 中，$\angle A:\angle B:\angle C=2:3:4$，求 $\angle D$。

	\textbf{\textcolor{CExample}{解题步骤：}}
	\begin{description}[style=nextline, leftmargin=2.8em, labelsep=0.5em, itemsep=0.45em]
		\item[\textbf{第一步（设未知数）：}] 设 $\angle A=2x$，$\angle B=3x$，$\angle C=4x$（$x>0$）。
			\par\textbf{理由：} 利用已知比例引入参数。
		\item[\textbf{第二步（列互补方程）：}] 由 $\angle A$ 与 $\angle C$ 互补，得 $2x+4x=180^\circ$，解得 $x=30^\circ$。
			\par\textbf{理由：} 圆内接四边形对角互补。
		\item[\textbf{第三步（求 $\angle D$）：}]
			\[
				\begin{aligned}
					\angle D &= 180^\circ - \angle B \\
					&= 180^\circ - 3\times 30^\circ \\
					&= 90^\circ .
			\end{aligned}
			\]
			（或利用四边形内角和为 $360^\circ$ 验算）。
			\par\textbf{理由：} $\angle B$ 与 $\angle D$ 互补。
	\end{description}

	\textbf{\textcolor{CExample}{关键结论：}}
	$\angle D=90^\circ$。

	\medskip
	\textbf{\textcolor{CExample}{例 3（综合应用）}}

	\textbf{\textcolor{CExample}{题目：}}
	在圆内接四边形 $ABCD$ 中，$\angle A$ 是 $\angle C$ 的 $2$ 倍，$\angle B$ 是 $\angle D$ 的 $3$ 倍，求四边形各内角的度数。

	\textbf{\textcolor{CExample}{解题步骤：}}
	\begin{description}[style=nextline, leftmargin=2.8em, labelsep=0.5em, itemsep=0.45em]
		\item[\textbf{第一步（设角）：}] 设 $\angle C=\alpha$，则 $\angle A=2\alpha$；设 $\angle D=\beta$，则 $\angle B=3\beta$。
			\par\textbf{理由：} 根据倍数关系设未知量。
		\item[\textbf{第二步（列互补方程）：}] 由对角互补得：
			\[
				\begin{aligned}
					\angle A+\angle C &= 2\alpha+\alpha \\
					&= 180^\circ,
			\end{aligned}
			\]
			\[
				\begin{aligned}
					\angle B+\angle D &= 3\beta+\beta \\
					&= 180^\circ .
			\end{aligned}
			\]
			\par\textbf{理由：} 圆内接四边形两组对角分别互补。
		\item[\textbf{第三步（解方程）：}] 由 $3\alpha=180^\circ$ 得 $\alpha=60^\circ$；由 $4\beta=180^\circ$ 得 $\beta=45^\circ$。
			\par\textbf{理由：} 解一元一次方程。
		\item[\textbf{第四步（回代求角）：}]
			\[
				\begin{aligned}
					\angle A &= 2\alpha \\
					&= 120^\circ , \\
					\angle C &= \alpha \\
					&= 60^\circ , \\
					\angle B &= 3\beta \\
					&= 135^\circ , \\
					\angle D &= \beta \\
					&= 45^\circ .
			\end{aligned}
			\]
			\par\textbf{理由：} 代入已知关系。
	\end{description}

	\textbf{\textcolor{CExample}{关键结论：}}
	$\angle A=120^\circ$，$\angle B=135^\circ$，$\angle C=60^\circ$，$\angle D=45^\circ$。
\end{examplebox}
